\subsection{Semi-barrelled spaces}

PROPOSITION 1. --- Let E be a locally convex space. The following conditions are equivalent :

\begin{enumerate}
\def\labelenumi{(\roman{enumi})}
\item
  Let U be a subset of E which absorbs every bounded subset of E, and which is the intersection of a sequence of convex, balanced and closed neighbourhoods of 0 in E. Then U is a neighbourhood of 0 in E.
\item
  For every locally convex space \(\mathbf{F}\) , every bounded subset of \(\mathcal{L}_{\mathfrak{b}}(\mathbf{E};\mathbf{F})\) which is the union of a countable family of equicontinuous subsets, is equicontinuous.
\item
  In the strong dual \(E_{b}^{\prime}\) of E, every bounded subset which is the union of a countable family of equicontinuous subsets, is equicontinuous.
\end{enumerate}

It is clear that (iii) is a particular case of (ii).

\begin{enumerate}
\def\labelenumi{(\roman{enumi})}
\item
  \(\Rightarrow\) (ii): let H be a bounded subset of \(\mathcal{L}_b(E;F)\) , and let \((\mathrm{H}_n)\) be a sequence of equicontinuous subsets of \(\mathcal{L}_b(E;F)\) such that \(\mathrm{H} = \bigcup_{n}\mathrm{H}_{n}\) . Let V be a convex, balanced and closed neighbourhood of 0 in F. For every \(n\) , the set \(\mathrm{W}_n = \bigcap_{u\in \mathrm{H}_n}u^{-1}(\mathrm{V})\) is a convex, balanced and closed neighbourhood of 0 in E since \(\mathrm{H}_n\) is equicontinuous. The set \(\mathrm{W} = \bigcap_{u\in \mathrm{H}}u^{-1}(\mathrm{V})\) absorbs every bounded subset of E, since H is bounded in \(\mathcal{L}_b(E;F)\) (III, p.~22), and we have \(\mathrm{W} = \bigcap_{n}\mathrm{W}_{n}\) . If E satisfies (i), then the set W is a neighbourhood of 0 in E, hence H is equicontinuous.
\item
  \(\Rightarrow\) (i): let \((\mathrm{U}_n)\) be a sequence of convex, balanced and closed neighbourhoods of 0 in E. We assume that the set \(\mathbf{U} = \bigcap_{n}\mathbf{U}_{n}\) absorbs every bounded subset of E, hence that its polar \(\mathbf{U}^{\circ}\) is bounded in \(\mathbf{E}_b^\prime\) . Then the set \(\mathbf{B} = \bigcup_{n}\mathbf{U}_{n}^{\circ}\) is contained in \(\mathbf{U}^{\circ}\) , hence is bounded in \(\mathbf{E}_b^\prime\) . If E satisfies (iii), the set B is equicontinuous in \(\mathbf{E}'\) ; consequently, the polar \(\mathbf{B}^{\circ} = \bigcap_{n}(\mathbf{U}_{n}^{\circ})^{\circ} = \bigcap_{n}\mathbf{U}_{n} = \mathbf{U}\) of B in E is a neighbourhood of 0 in E.
\end{enumerate}

DEFINITION 1. --- A locally convex space E is said to be semi-barrelled if it satisfies the equivalent conditions of prop. 1.

Every barrelled space is semi-barrelled. This is also true for every bornological space (III, p.~22, prop. 10).

